

\section{Chapter Overview}
The chapter highlights a brief discussion about the problem statement and the rationale behind undertaking this project. It starts by providing a brief introduction to the university admission process concerning the post-GST university admission process, wherein there are several stages of verification necessary to accomplish the entire process of admission. The rationale behind this project is based on the current state of affairs regarding the process of university admission. Presently, the admission process is carried out either manually or semi-manually, leading to inefficiencies in data management, which can be erroneous and slow, and poor communication between students and various departments within the university.Finally, the chapter states the objectives of the project.

\section{Background}

\noindent Admission of a university is an indispensable activity in every educational activity and is as old as education itself. With more students applying for admission than in previous years, admission process for admissions officers has become increasingly time-consuming. More applications resulted in heavier paperwork and processing challenges. Every admission is routed through various faculty for evaluation, and the manual admission process caused difficulty in admission, and archiving. 

\vspace{0.3cm}

\noindent In order to speed up and simplify the admission process, the development of the Automated Final Admission System is to enact document redesign and an automated management program. The proposed ``Digital Admission Procedure Management for Post-GST Candidates'' would store, route, and retrieve prospective and current application documents. The development of science and technology has immensely contributed to the growth of computers and the internet which has adversely increased the use and need for an automated applicant's admission system. 

\vspace{0.3cm}

\noindent Universities, Colleges, Companies, organizations, etc. are seeking better ways to distribute and disseminate information throughout the world. In order to achieve this, they develop and rely on online applications, which help them input, update, process, and disseminate information. The internet has increased the use of online applications because data could be easily and readily accessed as well as documented. This project is targeted at allowing admission of applicants' decisions to be made faster and more efficiently. It looks at improving the present traditional system of the manual process (paper-based) which is manually documented. This type of project is a hot topic in the whole world to digitalize our education system.